\chapter{Methodology}

% ====================================================================
\section{Overview}
% ====================================================================

This chapter presents the methodology adopted for developing an
integrated healthcare insurance fraud detection system that combines machine
learning with cryptographically secure blockchain technology. The
proposed framework proceeds through eight sequential stages: raw data collection
and merging, ECIES-based encryption of sensitive provider and beneficiary
information, SHA-256 blockchain storage with Merkle tree integrity verification,
preprocessing and provider-level feature engineering, Optuna-driven
hyperparameter optimisation, ensemble machine learning classification with
10-fold stratified cross-validation, threshold optimisation over the
precision-recall curve, and final performance evaluation augmented by SHAP and
LIME explainability analyses. Each stage is motivated by specific limitations
identified in the prior literature and is grounded in reproducible, open-source
implementations using methods available before June 2024. Together, these stages
form a complete, end-to-end pipeline that addresses both the predictive accuracy
and the auditability requirements of a production-grade fraud detection system
for the Centers for Medicare and Medicaid Services (CMS) context.

% ====================================================================
\section{Background Study}
% ====================================================================

Healthcare fraud imposes a severe and growing burden on health systems worldwide.
The United States Federal Bureau of Investigation estimates that fraudulent
billing, unnecessary procedures, and identity theft collectively drain tens of
billions of dollars annually from Medicare and Medicaid alone, with the
Congressional Budget Office placing the share of improper payments in the United
States public insurance system between 7 and 10 per cent of total expenditure
\cite{fbi2021,cbo2022}. Globally, the World Health Organisation has estimated
that between 10 and 25 per cent of all healthcare spending is lost to fraud,
waste, and abuse, with low- and middle-income countries bearing a
disproportionate share of this burden due to weaker audit infrastructure
\cite{who2010}. Fraudulent schemes range from upcoding and phantom billing by
providers to beneficiary identity theft and kickback arrangements between
suppliers and clinicians. The Medicare programme administered by the CMS is a
particularly attractive target because of its high claim volume, the complexity
of its reimbursement rules, and the limited capacity for manual pre-payment
auditing at scale.

Traditional fraud detection in healthcare has relied on retrospective rule-based
auditing, statistical outlier detection, and manual case review by trained
investigators. While these approaches establish a necessary baseline, they are
inadequate against sophisticated and adaptive fraud schemes that deliberately
stay below detection thresholds or exploit new billing codes. Machine learning
offers a data-driven alternative that can identify complex, non-linear patterns
across high-dimensional feature spaces derived from claims data, beneficiary
demographics, and provider behaviour. Blockchain technology complements ML by
providing a tamper-evident, append-only ledger for storing prediction records,
thereby enabling post-hoc auditability and regulatory compliance. The combination
of these two paradigms has attracted growing research interest since approximately
2018, with early work focusing on federated learning for privacy preservation and
more recent contributions exploring zero-knowledge proofs and smart-contract-based
audit trails. The present study operationalises these ideas on the publicly
available CMS-derived Kaggle dataset to produce a reproducible, benchmarked
reference implementation.

% ====================================================================
\section{Research Gap}
% ====================================================================

A systematic review of the existing literature on healthcare fraud detection
reveals that, while considerable progress has been made in applying individual
machine learning algorithms to insurance claims data, several important gaps
persist. These gaps concern the depth of feature engineering, the choice of
evaluation metrics, the handling of class imbalance, the integration of
explainability tools, the incorporation of cryptographic privacy guarantees, and
the end-to-end validation of blockchain-backed audit trails. The following
subsections characterise each gap and explain how the present study addresses it.

\subsection{Limited Utilisation of Advanced Technologies}

The majority of published healthcare fraud detection studies rely on a narrow set
of classical algorithms such as logistic regression, naive Bayes, and shallow
decision trees, without exploring the performance improvements available from
gradient-boosted ensembles or systematic hyperparameter search. Recent advances in
gradient boosting frameworks such as XGBoost \cite{chen2016}, LightGBM
\cite{ke2017}, and CatBoost \cite{prokhorenkova2018}, combined with automated
hyperparameter optimisation libraries such as Optuna \cite{akiba2019}, offer
substantially higher predictive accuracy on tabular imbalanced datasets. The
present study addresses this gap by benchmarking five classifiers and applying
Optuna-based Bayesian hyperparameter search with 60 trials per model, thereby
applying the state of the art in tabular machine learning within a rigorous
cross-validation protocol.

\subsection{Lack of Comprehensive Data Integration}

Many prior studies analyse only inpatient or outpatient claims in isolation,
forgoing the complementary diagnostic signals available when both data sources are
merged with beneficiary demographic records. This fragmented approach results in
incomplete provider-level profiles and diminishes the discriminative power of the
feature set. The present study addresses this gap by merging four raw CSV files
from the Kaggle rohitrox dataset, producing a unified claim-level table of
558,211 records that is subsequently aggregated to 5,410 provider-level
observation vectors with 190 engineered features, thereby capturing the full
breadth of each provider's billing behaviour and patient population.

\subsection{Dynamic Nature of Fraud Schemes}

Fraudulent providers frequently adapt their billing patterns in response to
detection systems, rendering static rule-based detectors obsolete within months
of deployment. A purely accuracy-driven ML model that maximises a fixed
probability threshold is similarly brittle when the distribution of fraud patterns
shifts. The present study addresses this challenge by supplementing standard
accuracy metrics with area under the precision-recall curve (AUC-PR), Matthews
Correlation Coefficient (MCC), and an explicit threshold optimisation step that
selects the operating point maximising the F1-score on held-out data, providing
a more robust and adjustable detection boundary than a fixed 0.5 cutoff.

\subsection{Blockchain Integration}

The existing fraud detection literature largely treats the prediction model as a
standalone artefact, providing no mechanism for immutably recording which
providers were flagged, when, by which model version, or under which
hyperparameter configuration. This absence of an auditable chain of custody
complicates regulatory compliance and post-hoc dispute resolution. The present
study addresses this gap by implementing a SHA-256 blockchain with Merkle tree
integrity verification that stores every fraud prediction record in tamper-evident
blocks, enabling independent verification of the audit trail without requiring
access to the raw data or model weights.

\subsection{Ethical and Privacy Concerns}

Healthcare data contains highly sensitive beneficiary information whose exposure
violates both HIPAA and ethical research norms. Studies that store raw beneficiary
identifiers alongside prediction outputs on shared ledgers or public repositories
introduce unacceptable re-identification risks. The present study addresses this
concern by encrypting all personally identifiable fields using the Elliptic Curve
Integrated Encryption Scheme (ECIES) with the secp256k1 curve, HKDF-SHA256 key
derivation, and AES-256-GCM authenticated encryption before any record is written
to the blockchain, ensuring that only holders of the authorised private key can
recover plaintext beneficiary information.

\subsection{Cost-Effectiveness}

Healthcare insurers and government audit agencies operate under strict budget
constraints that limit the number of providers that can be subjected to full
manual investigation. A fraud detection system that achieves high recall but very
low precision wastes investigative resources on false positives, while a system
that optimises precision at the expense of recall allows a large proportion of
fraudulent providers to escape detection. The present study addresses this
trade-off explicitly by reporting the full precision-recall curve, selecting an
operating threshold that balances precision and recall for the specific cost
structure of Medicare audits, and computing the F2-score in supplementary
analyses to reflect the higher societal cost of missed fraud compared to false
alarms.

\subsection{Regulatory Compliance}

Fraud detection systems deployed within the CMS ecosystem must satisfy audit
trail, data retention, and explainability requirements imposed by the Office of
Inspector General and the False Claims Act. Prior ML-based studies rarely address
these operational constraints, producing systems that perform well on benchmark
datasets but cannot be integrated into existing compliance workflows. The present
study addresses this gap by designing the blockchain storage layer to preserve
immutable, time-stamped records of every prediction event and by generating SHAP
summary plots and LIME explanations for every flagged provider, providing the
human-readable justifications required under emerging AI accountability
frameworks.

\subsection{User Adoption and Training}

Even technically superior fraud detection systems fail in practice if
investigators and auditors cannot understand or trust the system's outputs. Black-
box models that produce probability scores without feature-level justifications
are routinely rejected or underutilised by clinical and administrative staff.
The present study addresses this adoption barrier by integrating post-hoc
explainability at two levels: global feature importance via SHAP beeswarm and
bar plots that characterise population-level fraud drivers, and local instance
explanations via LIME that provide per-provider rationales directly suitable for
inclusion in audit reports submitted to human reviewers.

\subsection{Real-World Implementation Challenges}

Transitioning a fraud detection model from a research notebook to a production
system introduces challenges around reproducibility, containerisation, and
end-to-end data flow. Many published studies omit the engineering artefacts
necessary for real-world deployment. The present study addresses this gap by
providing a fully containerised pipeline implemented with Docker and
\texttt{docker-compose}, with each phase encapsulated in a separate Python script
that accepts command-line arguments for data directories, hyperparameter budgets,
and blockchain difficulty settings, ensuring that the results are fully
reproducible from the raw Kaggle CSV files.

\subsection{Benchmarking and Evaluation Metrics}

A recurring weakness in the healthcare fraud detection literature is the reliance
on accuracy or ROC-AUC as the primary evaluation metric. On highly imbalanced
datasets such as the CMS provider labels used here, where only 9.4 per cent of
providers are fraudulent, a naive classifier that predicts ``non-fraud'' for all
providers achieves 90.6 per cent accuracy. The present study addresses this by
reporting a full suite of metrics including F1-score, precision, recall,
AUC-PR, MCC, and ROC-AUC across 10-fold stratified cross-validation, providing
statistically meaningful comparisons between all five classifiers and their
SMOTE-ENN ablations.

% ====================================================================
\section{Dataset Description}
% ====================================================================

The dataset used in this study is the Healthcare Provider Fraud Detection Analysis
dataset published by \textit{rohitrox} on Kaggle, derived from publicly available
Centers for Medicare and Medicaid Services (CMS) Medicare claims records. The
dataset comprises four CSV files: \texttt{Train.csv} containing labels for 5,410
providers (``Yes'' or ``No'' for \texttt{PotentialFraud}),
\texttt{Train\_Beneficiarydata} with demographic and chronic condition records for
138,556 unique beneficiaries, \texttt{Train\_Inpatientdata} with 40,474 inpatient
hospitalisation claims, and \texttt{Train\_Outpatientdata} with 517,737 outpatient
visit claims. After merging and outer-joining inpatient and outpatient claims, then
joining with beneficiary demographics and provider labels, the resulting
claim-level table contains 558,211 rows. This merged table is subsequently
aggregated to 5,410 provider-level observation vectors with 190 engineered
features. The target variable is severely imbalanced: 506 providers (9.4\%) are
labelled as potential fraud and 4,904 (90.6\%) are labelled as non-fraud. Table
\ref{tab:dataset_features} describes the principal raw columns used in the
analysis.

\begin{table}[H]
\centering
\caption{Key columns in the Kaggle rohitrox Healthcare Provider Fraud Detection dataset}
\label{tab:dataset_features}
\begin{tabular}{p{4.5cm} p{9.5cm}}
\toprule
\textbf{Column Name} & \textbf{Description} \\
\midrule
\texttt{Provider} & Unique alphanumeric identifier for each healthcare provider (primary key for aggregation) \\
\texttt{BeneID} & Unique identifier for each Medicare beneficiary \\
\texttt{ClaimID} & Unique identifier for each individual insurance claim \\
\texttt{ClaimStartDt} & Date on which the insurance claim period began; used to compute claim duration and beneficiary age at claim time \\
\texttt{ClaimEndDt} & Date on which the claim period ended; difference from \texttt{ClaimStartDt} gives claim duration in days \\
{\footnotesize\texttt{InscClaimAmtReimbursed}} & Total amount reimbursed by Medicare for the claim in US dollars \\
\texttt{DeductibleAmtPaid} & Out-of-pocket deductible amount paid by the beneficiary \\
\texttt{AdmissionDt} & Date of hospital admission (inpatient claims only) \\
\texttt{DischargeDt} & Date of hospital discharge (inpatient claims only); used to compute length of stay \\
\texttt{DOB} & Beneficiary date of birth; used to compute patient age at claim time \\
\texttt{DOD} & Beneficiary date of death (if applicable); used to compute the dead-patient ratio per provider \\
\texttt{ChronicCond\_*} & Binary indicators (1 = yes, 2 = no) for 11 chronic conditions including Alzheimer's disease, heart failure, chronic kidney disease, cancer, obstructive pulmonary disease, depression, diabetes, ischaemic heart disease, osteoporosis, rheumatoid arthritis, and stroke \\
\texttt{AttendingPhysician} & Anonymised identifier of the attending physician; used to count unique physicians per provider \\
\texttt{ClmDiagnosisCode\_1} & Primary ICD diagnosis code for the claim; used to count unique diagnosis codes per provider \\
\texttt{PotentialFraud} & Binary target variable (``Yes''/``No'') indicating whether the provider has been identified as a potential fraudster \\
\bottomrule
\end{tabular}
\end{table}

% ====================================================================
\section{Research Methodology}
% ====================================================================

\subsection{Technique Used}

The proposed system integrates six complementary technical components to achieve
high-accuracy, explainable, and cryptographically auditable healthcare fraud
detection. These components are: (i) Optuna Bayesian hyperparameter optimisation
with a Tree-structured Parzen Estimator sampler, (ii) SMOTE-ENN hybrid sampling
for class imbalance analysis and ablation, (iii) the Elliptic Curve Integrated
Encryption Scheme for PII protection, (iv) five machine learning classifiers
encompassing gradient-boosted trees and ensemble methods, (v) a SHA-256 blockchain
with Merkle tree integrity verification for audit trail storage, and (vi) the
proposed eight-step end-to-end pipeline that orchestrates all of the above
components. Each component is described in detail below.

\begin{enumerate}[label=\roman*.]

% ---- i. Optuna ----
\item \textbf{Optuna Hyperparameter Optimisation.}
Optuna is a define-by-run automatic hyperparameter optimisation framework that
enables efficient search over high-dimensional, mixed-type parameter spaces
\cite{akiba2019}. In the present study, Optuna is configured with the
Tree-structured Parzen Estimator (TPE) sampler, which constructs two density
estimates - one over hyperparameter configurations that yield high objective
values and one over configurations that yield low objective values - and samples
candidates that maximise the ratio of these densities. For each of the five
classifiers, 60 Optuna trials are executed with a timeout of 1{,}800 seconds per
model. The inner-loop objective function is the area under the precision-recall
curve (AUC-PR) evaluated via 5-fold stratified cross-validation on the training
partition, chosen in preference to ROC-AUC because AUC-PR is more informative
under severe class imbalance. Each trial also includes a boolean hyperparameter
controlling whether SMOTE-ENN oversampling is applied within the cross-validation
pipeline, enabling the optimiser to discover the superior configuration for each
classifier without requiring the researcher to pre-commit to a resampling
strategy.

% ---- ii. SMOTE-ENN ----
\item \textbf{SMOTE-ENN for Class Imbalance.}
SMOTE-ENN is a hybrid resampling algorithm that combines the Synthetic Minority
Over-sampling Technique (SMOTE) with the Edited Nearest Neighbours (ENN) cleaning
rule \cite{batista2004}. SMOTE generates synthetic minority-class samples by
interpolating between existing minority instances in feature space, addressing
the under-representation of fraudulent providers. ENN subsequently removes
majority-class instances whose class label disagrees with the majority vote of
their three nearest neighbours, thereby cleaning borderline and noisy samples
from both classes and reducing class overlap. In the present pipeline,
SMOTE-ENN is applied inside an \texttt{imblearn} pipeline so that the resampling
occurs only on the training fold during cross-validation, preventing data leakage
from synthetic samples into the validation fold. A central finding of the study
is that, contrary to common assumptions, SMOTE-ENN consistently degrades
performance by 6 to 10 percentage points in F1-score relative to the baseline
cost-sensitive approach on this dataset, demonstrating that the
\texttt{scale\_pos\_weight} mechanism native to gradient-boosted classifiers is
a superior strategy for handling the 9.4 per cent fraud prevalence observed here.

% ---- iii. ECIES ----
\item \textbf{Elliptic Curve Integrated Encryption Scheme (ECIES).}
ECIES is a hybrid public-key encryption scheme that combines asymmetric key
agreement with symmetric authenticated encryption to provide both
confidentiality and integrity for arbitrarily large payloads \cite{shoup2001}.
The implementation in this study uses the secp256k1 elliptic curve for key
generation and the Elliptic Curve Diffie-Hellman (ECDH) primitive for shared-
secret derivation. The shared secret is processed through HKDF-SHA256 to derive
a 256-bit symmetric key, which is then used with AES-256 in Galois/Counter Mode
(AES-256-GCM) to encrypt the sensitive payload under a 96-bit random nonce. The
ciphertext transmitted to the blockchain record consists of the sender's ephemeral
compressed public key (33 bytes), the 12-byte nonce, the AES-GCM ciphertext, and
the 16-byte authentication tag, concatenated and Base64-encoded. This construction
ensures that provider PII such as beneficiary identifiers, diagnosis codes, and
reimbursement amounts are rendered unreadable to any party that does not hold
the designated private key, while the associated fraud score and Merkle hash
remain visible for external audit purposes.

% ---- iv. ML Classifiers ----
\item \textbf{Machine Learning Classifiers.}

Five supervised classifiers are evaluated in this study. Each is trained on the
provider-level feature matrix of 5,410 observations with 190 features, using
10-fold stratified cross-validation with Optuna-tuned hyperparameters.

\textbf{XGBoost (eXtreme Gradient Boosting)} is a scalable, regularised gradient
boosting framework that constructs an additive ensemble of decision trees by
greedily minimising a differentiable loss function augmented with $L_1$ and $L_2$
regularisation terms \cite{chen2016}. XGBoost supports native handling of class
imbalance through the \texttt{scale\_pos\_weight} parameter, which up-weights the
minority class in the loss function by a factor proportional to the class-
frequency ratio. Tuned hyperparameters optimised in this study include the number
of estimators (100 to 800), maximum tree depth (3 to 10), learning rate (0.01 to 0.3,
log-uniform), subsample ratio, column subsampling ratio, minimum child weight,
gamma regularisation, and the $L_1$ and $L_2$ regularisation coefficients.

\textbf{LightGBM (Light Gradient Boosting Machine)} improves upon conventional
gradient boosting through two key innovations: Gradient-based One-Side Sampling
(GOSS), which retains instances with large gradients while randomly sampling
those with small gradients to reduce training time, and Exclusive Feature Bundling
(EFB), which merges mutually exclusive sparse features to reduce
dimensionality \cite{ke2017}. LightGBM grows trees leaf-wise rather than
level-wise, enabling deeper, more specialised splits with fewer iterations.
Its native \texttt{is\_unbalance} and \texttt{scale\_pos\_weight} options provide
flexible imbalance correction. In the Optuna search, the number of leaves, minimum
child samples, feature fraction, bagging fraction, and regularisation terms are
tuned alongside the core learning rate and estimator count.

\textbf{CatBoost} is a gradient boosting library developed by Yandex that
introduces ordered boosting to eliminate prediction shift caused by target
leakage in categorical feature encoding \cite{prokhorenkova2018}. Rather than
label-encoding categorical columns prior to training, CatBoost applies an
online permutation-based statistic that prevents the model from seeing the
target value of a sample during its own gradient computation. Although the
provider-level features in this study are predominantly numeric after aggregation,
CatBoost's symmetric tree structure and built-in overfitting detector make it a
competitive baseline for the size of the dataset. The class imbalance is handled
via the \texttt{auto\_class\_weights} parameter (with Balanced, SqrtBalanced, and
None as candidate values), and the depth, learning rate, $L_2$ regularisation, and
number of iterations are included in the Optuna search space.

\textbf{Gradient Boosting Classifier} (scikit-learn implementation) applies the
classical Friedman gradient boosting algorithm \cite{friedman2001}, fitting a
fixed-depth regression tree to the negative gradient of the loss function at each
stage. Unlike the later XGBoost and LightGBM variants, the scikit-learn
implementation lacks native GPU acceleration and sophisticated leaf-pruning
heuristics, but provides a well-validated baseline with straightforward
hyperparameter interpretation. The class imbalance is addressed by passing
\texttt{sample\_weight} arrays with weights proportional to the inverse class
frequency. Tuned parameters include the number of estimators, maximum depth,
learning rate, minimum samples per split, and the subsampling fraction.

\textbf{Random Forest} is a bagging ensemble that trains a large number of
decision trees on bootstrap replicates of the training data, selecting a random
subset of features at each split node to decorrelate the individual trees
\cite{breiman2001}. Prediction is made by majority vote (classification) or
averaging (regression). Random Forest is particularly well suited to datasets
with moderate sample sizes and many features because the bagging variance
reduction complements the feature-subsampling bias reduction. Class imbalance is
addressed by setting \texttt{class\_weight=`balanced\_subsample'}, which reweights
samples by the inverse of their class frequency within each bootstrap subsample.
The number of trees, maximum depth, minimum samples per leaf, and the feature
subsampling fraction are included in the Optuna search.

% Logistic Regression was evaluated in early experiments but excluded from the
% final pipeline due to consistently lower performance on this dataset compared
% to the five gradient-boosted and ensemble-based classifiers retained above.

\paragraph{Cross-Validation Feature Computation.}
All cross-provider features (z-score normalisations, percentile ranks, physician
billing statistics, and network metrics) are computed within each
cross-validation fold using only training-fold data, preventing information
leakage from test providers.

% ---- v. Blockchain ----
\item \textbf{SHA-256 Blockchain with Merkle Trees.}
The blockchain module implemented in this study uses SHA-256 cryptographic
hashing, Merkle trees, and SQLite-backed persistence. Each
block stores a batch of up to 50 fraud prediction records and records the
SHA-256 hash of the concatenation of its own block index, timestamp, batch
Merkle root, and the hash of the previous block, thereby forming a tamper-evident
chain in which any alteration of a historical record invalidates all subsequent
block hashes. The Merkle tree is computed bottom-up by hashing each record
dictionary as a JSON string at the leaf level, then pairwise-hashing adjacent
leaves until a single root digest is obtained. The Merkle root enables efficient
inclusion proofs: to verify that a specific prediction record belongs to a given
block, an auditor need only follow the $O(\log n)$ sibling hashes in the proof
path rather than re-hashing all records in the block. Proof-of-work with
adjustable difficulty (configurable leading zero count) further hardens the chain
against retrospective modification. Chain validation traverses all blocks from
genesis to the current tip, recomputing each block hash and Merkle root and
verifying that hash linkage and proof-of-work constraints are satisfied, providing
a complete integrity certificate for the audit trail.

\paragraph{Design Rationale: Custom Blockchain.}
The custom Python blockchain is chosen over enterprise frameworks such as Hyperledger Fabric and Ethereum for three reasons. First, \emph{transparency}: every cryptographic operation --- SHA-256 hashing, Merkle proof construction, ECIES encryption --- is implemented in the source code and can be inspected, modified, and explained at the implementation level. Second, \emph{reproducibility}: the system runs on any machine with Python and SQLite, requiring no Docker infrastructure, Go chaincode compilation, or certificate authority configuration. Third, \emph{research focus}: the thesis contribution is the integration of blockchain audit trails with an ML fraud detection pipeline, not the blockchain infrastructure itself. A production deployment would use Hyperledger Fabric for multi-organisation consensus and enterprise-grade access control, as discussed in Section~5.4 (Future Work).

\paragraph{Proof-of-Work Difficulty.}
The PoW difficulty parameter is set to 2 (block hashes must begin with two leading zero hex digits), corresponding to approximately 256 hash attempts on average. This level is calibrated for audit-trail integrity verification rather than adversarial mining resistance. The complete 110-block chain builds in 7.0 seconds (1.3\,ms per record), well within real-time audit requirements. Higher difficulty values provide no additional security benefit in a single-node research context; a distributed deployment would adopt committee-based consensus (e.g., PBFT or Raft) instead.

\paragraph{Threat Model.}
The blockchain component addresses three specific threats:
\begin{enumerate}[nosep]
\item \textbf{Post-hoc record tampering}: an unauthorised party modifies a stored fraud score after the fact. Mitigated by SHA-256 hash chaining and Merkle proofs, which make any modification detectable through chain revalidation.
\item \textbf{PII exposure}: an unauthorised party reads provider identity from the audit trail. Mitigated by ECIES field-level encryption; only holders of the designated private key can recover plaintext provider identifiers.
\item \textbf{Repudiation}: an investigator denies having flagged a provider. Mitigated by timestamped, hash-linked audit entries that provide a non-repudiable record of every prediction event.
\end{enumerate}
Byzantine fault tolerance, Sybil attacks, and 51\% mining attacks are explicitly out of scope for this single-node research prototype and are listed as future work in Section~5.4.

\end{enumerate}

\subsubsection{Design Decisions and Rationale}
Table~\ref{tab:design} summarises the key design decisions, the alternatives considered, and the rationale for each choice. Every selection is motivated by either empirical evidence from the experiments, alignment with best practices in the literature, or practical constraints of the research setting.

\begin{longtable}{p{2.8cm} p{3cm} p{2.8cm} p{5cm}}
\caption{Design decisions and rationale for key architectural choices.} \label{tab:design} \\
\toprule
\textbf{Category} & \textbf{Choice} & \textbf{Alternatives} & \textbf{Rationale} \\
\midrule
\endfirsthead
\multicolumn{4}{l}{\textit{Table~\ref{tab:design} continued}} \\
\toprule
\textbf{Category} & \textbf{Choice} & \textbf{Alternatives} & \textbf{Rationale} \\
\midrule
\endhead
\bottomrule
\endfoot
Language & Python 3.10 & R, Java, Scala & Richest ML ecosystem (scikit-learn, XGBoost, SHAP); native integration with blockchain module \\
Dataset & Kaggle rohitrox & Real CMS, synthetic & Publicly available; peer-benchmarkable; real Medicare claims structure \\
Aggregation & Provider-level & Claim-level & Fraud labels exist at provider level; claim-level requires label propagation \\
Features & 190 engineered & 52 core subset & Ablation study (Table~4.6) confirms full set outperforms all subsets \\
Classifiers & 5 tree-based & SVM, KNN, MLP & Gradient boosting dominates tabular benchmarks; LogReg excluded after Optuna showed lowest AUC-PR \\
Ensemble & Optuna-weighted & Stacking, average & AUC-PR-optimised weights; stacking overfits with only 5 base models \\
HPO framework & Optuna TPE & Grid, random search & TPE converges faster on mixed-type spaces (60 trials per model) \\
CV strategy & 10-fold stratified & 5-fold, holdout & 10 folds yield 10 paired samples for Wilcoxon; stratification preserves 9.4\% fraud rate \\
Imbalance & Native class weighting & SMOTE-ENN, ADASYN & Empirical: SMOTE degrades F1 by 7--10\% (Table~4.5) \\
Primary metric & F1-score & Accuracy, AUC-ROC & Accuracy misleading at 90.6\% class prior; F1 balances precision and recall \\
Threshold & $t = 0.444$ & $t = 0.5$ & Validation-optimal; fixed for all downstream evaluation \\
Blockchain & Custom SHA-256 & Hyperledger, Ethereum & Cryptographic transparency; no Docker/Go dependency; reproducible \\
Hash function & SHA-256 & SHA-3, Blake2 & NIST standard; Bitcoin/TLS compatibility; widest library support \\
Encryption & ECIES (secp256k1) & RSA, secp256r1 & Bitcoin-compatible curve; authenticated encryption without separate signing \\
KDF & HKDF-SHA256 & PBKDF2, scrypt & HKDF designed for ECDH shared secrets (RFC~5869); PBKDF2 for passwords \\
Symmetric cipher & AES-256-GCM & AES-CBC, ChaCha20 & GCM provides authenticated encryption in single pass; CBC needs separate HMAC \\
Merkle tree & Binary SHA-256 & Patricia, Verkle & Simplest; sufficient for batch verification; Patricia for key-value state only \\
PoW difficulty & 2 & Higher & Audit-trail integrity, not adversarial resistance; 7\,s for 110 blocks \\
Block size & 50 records & Dynamic & Balances Merkle depth (6 levels) with mining frequency \\
Persistence & SQLite & PostgreSQL, LevelDB & Zero-configuration, ACID-compliant, single-file portability \\
Web framework & Flask & Django, FastAPI & Lightweight Jinja2 templates; Django overhead unnecessary \\
Explainability & SHAP + LIME & SHAP only & Dual methods cross-validate attributions (global + local) \\
Statistical tests & Friedman + Wilcoxon & $t$-test, ANOVA & Non-parametric; no normality assumption on fold F1 distributions \\
Bootstrap & 2000 iterations & 1000, 5000 & Stable 95\% CI; diminishing returns beyond 2000 \\
Feature selection & All 190 & Top-$K$, PCA, RFE & Full set outperforms subsets; trees handle redundancy via subsampling \\
\end{longtable}

% ---- vi. Proposed Methodology / Flowchart ----
\subsubsection{Pipeline Flowchart}

The proposed eight-step end-to-end pipeline is illustrated in Figure
\ref{fig:pipeline_flowchart}. The pipeline proceeds sequentially from raw data
ingestion through to blockchain-backed prediction storage, with the Optuna
hyperparameter optimisation loop embedded between the preprocessing and training
stages.

\begin{figure}[H]
\centering
\begin{tikzpicture}[
    node distance=0.8cm and 0.5cm,
    box/.style={
      rectangle, rounded corners=4pt, draw=black, thick,
      minimum width=5.2cm, minimum height=0.7cm,
      text centered, font=\small\bfseries, fill=gray!12
    },
    arrow/.style={->, thick, >=stealth, color=gray!70}
  ]
  % Row 1 (left to right)
  \node[box, fill=green!12] (step1) {Step 1: Data Collection};
  \node[box, fill=green!12, right=of step1] (step2) {Step 2: ECIES Encryption};

  % Row 2 (right to left)
  \node[box, fill=red!8, below=of step2] (step3) {Step 3: Blockchain Storage};
  \node[box, fill=blue!10, below=of step1] (step4) {Step 4: Preprocessing};

  % Row 3 (left to right)
  \node[box, fill=orange!12, below=of step4] (step5) {Step 5: Optuna HPO};
  \node[box, fill=orange!12, right=of step5] (step6) {Step 6: ML Training (10-CV)};

  % Row 4 (right to left)
  \node[box, fill=yellow!15, below=of step6] (step7) {Step 7: Threshold Opt.};
  \node[box, fill=yellow!15, below=of step5] (step8) {Step 8: Evaluation + SHAP};

  % Snake arrows
  \draw[arrow] (step1) -- (step2);
  \draw[arrow] (step2) -- (step3);
  \draw[arrow] (step3) -- (step4);
  \draw[arrow] (step4) -- (step5);
  \draw[arrow] (step5) -- (step6);
  \draw[arrow] (step6) -- (step7);
  \draw[arrow] (step7) -- (step8);

\end{tikzpicture}
\caption{End-to-end eight-step pipeline for ML-based healthcare fraud detection
with ECIES encryption and SHA-256 blockchain audit trail. The snake layout shows
the sequential flow through data ingestion and encryption (green), blockchain
storage (red), preprocessing (blue), model training and tuning (orange), and
evaluation (yellow) stages.}
\label{fig:pipeline_flowchart}
\end{figure}

\noindent The following subsections describe each step in detail.

\vspace{0.5em}
\noindent\textbf{Step 1: Data Collection.}

\begin{enumerate}[label=\roman*.]
  \item The four raw CSV files from the Kaggle rohitrox Healthcare Provider Fraud
  Detection Analysis dataset are downloaded and verified. The files are
  \texttt{Train.csv} (5,410 provider-level fraud labels),
  \texttt{Train\_Beneficiarydata} (138,556 beneficiary demographic records),
  \texttt{Train\_Inpatientdata} (40,474 inpatient claims), and
  \texttt{Train\_Outpatientdata} (517,737 outpatient claims).
  \item The inpatient and outpatient tables are outer-joined on their common
  columns to produce a unified claims table of 558,211 records, with an
  \texttt{is\_inpatient} indicator column distinguishing the two claim types.
  \item The claims table is then left-joined with the beneficiary demographic
  table on \texttt{BeneID} and subsequently with the provider fraud labels on
  \texttt{Provider}, yielding the complete claim-level dataset from which all
  downstream features are derived.
\end{enumerate}

\vspace{0.5em}
\noindent\textbf{Step 2: ECIES Encryption.}

\begin{enumerate}[label=\roman*.]
  \item A secp256k1 key pair is generated for the audit authority. The public key
  is used to encrypt sensitive fields; the private key is held offline by the
  audit authority and never stored in the pipeline artefacts.
  \item For each provider record scheduled for blockchain storage, the personally
  identifiable fields - including beneficiary identifiers, attending physician
  codes, and raw reimbursement amounts - are serialised to JSON and encrypted
  under the audit authority's public key using the ECIES scheme described in
  Section 3.5(iii).
  \item The encrypted ciphertext blob and its HKDF-derived key material are stored
  in the blockchain record alongside the plaintext fraud score and provider
  pseudonym, ensuring that the audit log is fully queryable by fraud score while
  sensitive PII remains inaccessible without the private key.
\end{enumerate}

\vspace{0.5em}
\noindent\textbf{Step 3: SHA-256 Blockchain Storage.}

\begin{enumerate}[label=\roman*.]
  \item Provider fraud prediction records, each containing the provider
  pseudonym, fraud probability, binary prediction, and ECIES-encrypted PII,
  are batched into groups of up to 50 records per block.
  \item For each batch, a SHA-256 Merkle root is computed from the JSON
  serialisations of the individual records, and the block header is constructed
  by concatenating the block index, Unix timestamp, Merkle root, and previous
  block hash, then hashing this header with SHA-256 to produce the block hash.
  \item Proof-of-work mining iterates a nonce field until the block hash satisfies
  the configured difficulty constraint (a prescribed number of leading zero hex
  digits), and the completed block is appended to the SQLite-backed chain. Chain
  validity is confirmed after each insertion by re-traversing the full chain.
\end{enumerate}

\vspace{0.5em}
\noindent\textbf{Step 4: Preprocessing and Feature Engineering.}

\begin{enumerate}[label=\roman*.]
  \item Date columns (\texttt{ClaimStartDt}, \texttt{ClaimEndDt},
  \texttt{AdmissionDt}, \texttt{DischargeDt}, \texttt{DOB}, \texttt{DOD}) are
  parsed and used to compute derived numeric quantities: claim duration in days,
  inpatient length of stay, and beneficiary age at the time of each claim.
  Deceased beneficiary status is derived from the presence of a non-null
  \texttt{DOD} value.
  \item Provider-level aggregation computes, for each of the 5,410 providers,
  statistics including total claim count, unique beneficiary count, unique
  physician count, unique diagnosis code count, total and maximum reimbursement
  amounts, total deductible paid, mean claim duration, the proportion of inpatient
  claims, the mean and maximum patient age, the dead-patient ratio (deceased
  beneficiaries divided by total unique beneficiaries), and the mean count of
  chronic conditions per beneficiary.
  \item Infinite values introduced by division operations are replaced with
  \texttt{NaN} and subsequently imputed with column medians; all 190 resulting
  features are verified to contain no remaining missing values before further
  processing.
\end{enumerate}

\vspace{0.5em}
\noindent\textbf{Step 5: Optuna Hyperparameter Optimisation.}

\begin{enumerate}[label=\roman*.]
  \item For each of the five classifiers, an Optuna \texttt{Study} object is
  created with a \texttt{TPESampler} seeded at 42 and a
  \texttt{MedianPruner} that terminates unpromising trials early based on
  intermediate AUC-PR values reported at each inner cross-validation fold.
  \item The objective function evaluates each hyperparameter configuration via
  5-fold stratified cross-validation using the \texttt{average\_precision\_score}
  as the metric, with the complete pipeline - including optional SMOTE-ENN
  resampling controlled by a boolean trial parameter - encapsulated in an
  \texttt{imblearn} pipeline to prevent data leakage.
  \item After 60 trials or 1{,}800 seconds (whichever is reached first), the
  best trial's hyperparameters and the optimal SMOTE-ENN decision are serialised
  to \texttt{results/best\_params.json} for reproducible reuse in Phase 5.
\end{enumerate}

\vspace{0.5em}
\noindent\textbf{Step 6: ML Ensemble Classifier (10-fold CV).}

\begin{enumerate}[label=\roman*.]
  \item Each classifier is instantiated with the Optuna-tuned hyperparameters
  loaded from \texttt{best\_params.json} and evaluated under 10-fold stratified
  cross-validation, ensuring that the fraud-to-non-fraud ratio is maintained in
  every fold.
  \item The held-out fold predictions are aggregated to produce out-of-fold
  probability scores for all 5,410 providers, from which the macro-averaged F1,
  precision, recall, AUC-PR, MCC, and ROC-AUC are computed with 95\% confidence
  intervals derived from 1{,}000 bootstrap resamples.
  \item In addition to the five individual classifiers, an Optuna-weighted
  ensemble combining XGBoost, LightGBM, and CatBoost with optimised model
  weights is evaluated in the same cross-validation framework, quantifying
  the additional predictive gain achievable by combining the base classifiers.
\end{enumerate}

\vspace{0.5em}
\noindent\textbf{Step 7: Threshold Optimisation.}

\begin{enumerate}[label=\roman*.]
  \item The precision-recall curve is computed from the out-of-fold probability
  scores of the best-performing classifier (Optuna-tuned XGBoost) using the
  \texttt{precision\_recall\_curve} function from scikit-learn.
  \item For each candidate threshold value along the curve, the F1-score is
  computed as the harmonic mean of precision and recall at that operating point,
  and the threshold that maximises F1 is reported for reference.
  \item For the final binary predictions, a fixed classification threshold of
  $t = 0.444$ is applied to convert probability outputs, ensuring consistent
  evaluation across all folds without threshold overfitting.
\end{enumerate}

\vspace{0.5em}
\noindent\textbf{Step 8: Performance Evaluation and Explainability.}

\begin{enumerate}[label=\roman*.]
  \item The full evaluation suite reports F1-score, precision, recall, AUC-PR,
  MCC, and ROC-AUC for all five classifiers and their SMOTE-ENN ablations, with
  bootstrap confidence intervals, paired Wilcoxon signed-rank tests for
  statistical significance, and confusion matrices.
  \item Global feature importance is visualised using SHAP beeswarm and bar plots
  computed from the Optuna-tuned XGBoost model, identifying deductible payment
  patterns, maximum reimbursement amounts, and claim duration anomalies as the
  strongest fraud indicators at the population level.
  \item Local per-provider explanations are generated with LIME using a tabular
  explainer with 500 perturbation samples per instance, providing human-readable
  feature contribution summaries for each flagged provider that are suitable for
  inclusion in regulatory audit reports.
\end{enumerate}
